\section{Глава~\ref{sec:dofaalt}. Другие теории \nt{DOPA}}

Каждая из расширенных картин опирается на свои прямые измерения дофамина (разброс точек переворота, ответы на неприятное и новое, медленный рост, ответы на смену вкуса), но формулы, которые их объясняют, --- теории, и между двумя из них опыты пока противоречат друг другу.

{\footnotesize
\setlength{\tabcolsep}{3pt}\renewcommand{\arraystretch}{1.15}
\begin{longtable}{>{\raggedright}p{0.5cm}>{\raggedright}p{4.0cm}>{\raggedright}p{5.6cm}>{\raggedright}p{2.3cm}>{\raggedright\arraybackslash}p{1.7cm}}
\toprule
№ & Утверждение & Опыт и что измерено & Источник & Статус \\
\midrule\endfirsthead
\toprule
№ & Утверждение & Опыт и что измерено & Источник & Статус \\
\midrule\endhead
1 & Асимметричное правило~\eqref{eq:asymmetric} сходится к точке~\eqref{eq:expectile}; при $\kappa=1/2$ это среднее, для капли с вероятностью $1/2$ --- $V=\kappa$ (рис.~\ref{fig:expectiles}) & Точка~\eqref{eq:expectile} --- экспектиль, решение задачи наименьших квадратов с разными весами положительных и отрицательных отклонений; для примера главы вывод в самой главе. & \cite{NeweyPowell1987}; вывод в главе & теория \\
2 & У разных \nt{DOPA}-нейронов разные точки переворота, и они упорядочены по асимметрии (утверждение~\ref{prop:expectile}) & Мыши, оптогенетически помеченные \nt{DOPA}-нейроны вентральной покрышки, награды разного размера: точка, где всплеск сменяется провалом, у разных нейронов разная; у нейронов с большей асимметрией ответов она выше; по ответам группы восстанавливается форма распределения наград. Что это главная функция разброса --- гипотеза. & \cite{Dabney2020} & измерено \\
3 & Часть \nt{DOPA}-нейронов отвечает на неприятное всплеском & Обезьяны: одни \nt{DOPA}-нейроны возбуждаются на предвестник награды и тормозятся на предвестник струи воздуха в глаз, другие возбуждаются на оба; группы расположены в разных частях среднего мозга. & \cite{Matsumoto2009, BrombergMartin2010} & измерено \\
4 & Модуль ошибки или новизна задаёт скорость обучения & Прямого опыта, где сигнал важности у \nt{DOPA}-нейронов меняет скорость обучения, в цитируемых работах нет; обзор предлагает роль сигнала «внимание, важно». & \cite{BrombergMartin2010} & гипотеза \\
5 & Отдельный провод для оси опасности & Мыши: \nt{DOPA}-нейроны, идущие к заднему концу стриатума, отвечают на новые и сильные раздражители, а не на награду; их разрушение ослабляет избегание нового предмета. & \cite{Menegas2018} & измерено \\
6 & Оптимальное время действия $\tau_a^*=\sqrt{C/\bar R}$~\eqref{eq:vigor} & Минимум суммы $C/\tau_a+\bar R\tau_a$; вывод в главе. В исходной модели скорость действий задаёт цена времени, равная средней награде. & \cite{Niv2007}; вывод в главе & теория \\
7 & Медленный уровень дофамина несёт $\bar R$ и задаёт скорость действий (утверждение~\ref{prop:multiplex}) & Крысы, микродиализ и вольтамперометрия в прилежащем ядре: уровень дофамина от минуты к минуте следует за частотой наград и за скоростью действий. У людей L-DOPA усиливает зависимость времени реакции от средней частоты наград. Что медленный уровень равен именно $\bar R$, не доказано. & \cite{Hamid2016, Beierholm2013vigor} & гипотеза \\
8 & Медленный уровень складывается из двух источников: выброс меняется у выходов без изменения частоты импульсов, но медленный рост бывает и в самой частоте & Крысы, одна задача: выброс в прилежащем ядре следует за меняющимся ожиданием награды, а частота импульсов опознанных \nt{DOPA}-нейронов вентральной покрышки --- нет. Мыши в виртуальном коридоре (строка~10): плавный рост при приближении к награде есть и в импульсах опознанных \nt{DOPA}-нейронов. & \cite{Mohebi2019, Kim2020} & измерено \\
9 & Дофамин плавно растёт, пока животное приближается к далёкой награде & Крысы в лабиринте, вольтамперометрия в стриатуме: концентрация дофамина нарастает за секунды по мере приближения к цели, крутизна зависит от расстояния и размера награды. & \cite{Howe2013ramp, Hamid2016} & измерено \\
10 & Рост --- ошибка $\delta$ при уточнении оценки~\eqref{eq:ramp}; скачок в коридоре даёт всплеск (рис.~\ref{fig:ramp}) & Формула и число $\delta=0.75\,V$ в секунду --- вывод в главе. Опыт: мыши в виртуальном коридоре, мгновенный перенос ближе к цели; импульсы тел, кальций тел и выходов и концентрация в стриатуме отвечают как ошибка, а не как ожидание $V$. Какое из двух объяснений верно на всех выходах, обсуждается. & \cite{Kim2020} & измерено \\
11 & Взгляд назад: дофамин сообщает меру причинной связи~\eqref{eq:retro} & Мыши, выброс дофамина в прилежащем ядре в опытах, где эта теория и ошибка~\eqref{eq:ctd} предсказывают разное: в ряде опытов выброс следовал первой. & \cite{Jeong2022} & гипотеза \\
12 & Ответ дофамина при обучении постепенно сдвигается от награды к предвестнику & Мыши, сигнал выходов \nt{DOPA}-нейронов в брюшном стриатуме по ходу обучения: всплеск постепенно переходит с момента награды на момент предвестника, как требует обучение по~\eqref{eq:ctd}. & \cite{Amo2022} & измерено \\
13 & \nt{DOPA}-нейроны отвечают на смену вкуса награды при той же ценности & Крысы, запись нейронов вентральной покрышки: замена сока другим, равным по ценности, вызывает ответ, хотя ошибка~\eqref{eq:ctd} равна нулю. & \cite{Takahashi2017} & измерено \\
14 & Искусственные всплески учат связи двух нейтральных раздражителей & Крысы, опыт с предварительным сочетанием двух раздражителей без награды: оптогенетическое включение \nt{DOPA}-нейронов на переходе от первого ко второму создаёт связь, выключение мешает её выучить. & \cite{Sharpe2017} & измерено \\
15 & Вектор ошибок по признакам~\eqref{eq:vectortd} & Теория ошибки предсказания признаков; прямой проверки векторной формы нет. & \cite{Gardner2018} & гипотеза \\
16 & Разные \nt{DOPA}-нейроны в одной задаче несут разные величины & Мыши в виртуальном коридоре, двухфотонная съёмка \nt{DOPA}-нейронов: одни связаны с наградой, другие со скоростью и поворотом, третьи с предвестниками и точностью выбора. & \cite{Engelhard2019} & измерено \\
17 & Жгут вместо провода (утверждение~\ref{prop:bundle}) & Сводка строк~2--16: каждое разложение измерено по отдельности; единой картины, что все они --- части одного сигнала, опытом нет. & --- & гипотеза \\
\bottomrule
\end{longtable}}
